\clearpage
\section{Wyznaczenia masy startowej i masy własnej samolotu}
    
    \paragraph{}
        Z racji iz projektnie słozy  do obserwacji i nie rzuca ładubku w trakcie lotu,
        Do obliczeń masy ładunku samolotu wybrałem methode nr 2 opisana w \cite{galinski},
        Założeni i informacjęi najprawdopowdopniej do projektu, masa zastosowanej głowicy GOD-1 firmy PCO S,A zakładam 70 kg.
        
    \subsection{Obliczenie stosunków mas przed i po każdej fazie lotu}
        \begin{itemize}
            \item Rozruch:
                \begin{align}
                    \frac{W_{1}}{W_{0}}=0,97
                \end{align}
            \item Rozbieg:
                \begin{align}
                    \frac{W_{2}}{W_{1}}=0,985
                \end{align}
            \item Wznoszenie:
                \begin{align}
                    \frac{W_{3}}{W_{2}}=0,999
                \end{align}           
            \item Dolot:
                \begin{align}
                    \frac{W_{4}}{W_{3}}=0,984
                \end{align} 
            \item Wznoszenie:
                \begin{align}
                    \frac{W_{5}}{W_{4}}=0,997
                \end{align}  
            \item Patrolowanie:
                \begin{align}
                    \frac{W_{6}}{W_{5}}=0,778
                \end{align}
            \item Opadanie:
                \begin{align}
                    \frac{W_{7}}{W_{6}}=0,993
                \end{align}  
            \item Powrót:
                \begin{align}
                    \frac{W_{8}}{W_{7}}=0.993
                \end{align}  
            \item Opadanie:
                \begin{align}
                    \frac{W_{9}}{W_{8}}=0.994
                \end{align}  
            \item Oczekiwanie:
                \begin{align}
                    \frac{W_{10}}{W_{9}}=0.999
                \end{align}  
            \item Opadanie
                \begin{align}
                    \frac{W_{11}}{W_{10}}=0.994
                \end{align}
            \item Ladowanie:
                \begin{align}
                    \frac{W_{12}}{W_{11}}= 0.995
                \end{align}
        \end{itemize}

        \begin{align}
            \frac{W_{f}}{W_{0}} = 0.308
        \end{align}

    \subsection{Podsumowanie obliczenia massy staretowej i pustego samolotu}

        \paragraph{}
            Następnie na podstawie methody iteracjnej została wyznaczona nastepująca masa.
            masa obliczona 881 masa założona: 880.96 roznica około 1 promila wiec okk.